\section{Check completed routines through their public interfaces}
The observation recipes describe particular stops while a routine is being built.
A completed routine needs a different test: call it with inputs whose result is known,
then compare its externally visible state. The public checks in
\path{packages/030-qoi/checks} make those calls through the six published entry
points. They do not inspect the routine's labels, scratch registers or frame layout.
Keep the supplied caller, data arenas and \texttt{function\_return} label so that the
checks can supply inputs and regain control after each call.

Run \texttt{java -jar /path/to/cli.jar grade /path/to/check.kts /path/to/working.asm}
with one check file and the current working assembly source. The \texttt{bound},
\texttt{validate}, \texttt{decode} and \texttt{encode} check files exercise their named
interfaces. The \texttt{format} file supplies malformed streams and impossible byte
counts to the validator and decoder. The \texttt{files} file supplies in-memory files,
short transfers and failures. These file fixtures are installed by the grader; the
check does not open the student's real input or output files. The \texttt{caller} file
runs the connected program through \texttt{main} and checks its published result.

An encoder check decodes the returned stream with a host implementation of the local
format rules, then compares dimensions, descriptive flags and every RGBA component.
It also checks the stated size bound for uniform images. Returning a different legal
chunk sequence is therefore acceptable. A decoder check compares the complete output
words, including alpha, rather than judging the RGB display. Error cases check that
unpublished metadata and output remain untouched. Memory-access traces also check
whether a routine reads or writes outside its caller's permitted intervals, even when
an accidental access happens to leave the final bytes unchanged.

These are feedback checks for completed interfaces. An unfinished development return
of four is expected to fail them. Continue using the step's named checkpoint while
that routine is incomplete. Passing the public checks is evidence about their stated
cases; additional private checks exercise other inputs and repeated calls. No exact
encoding or instruction-count comparison is required for a free-writing solution.
